\chapter{Répertoire du vocabulaire religieux}\label{ch:rep_religieux}

Ce répertoire donne toutes les occurrences du vocabulaire religieux relevées dans le catalogue (ensemble A), classées par emploi puis par forme. Pour chaque occurrence, on indique la notice, l’année de classement, la forme développée et son statut textuel ; une note signale les cas où la confrontation des témoins ou la description matérielle modifie l’interprétation (monogramme probable, restitution divergente, lecture contestée). Les occurrences des notices exclues de l’ensemble C sont signalées. Le répertoire est produit par le programme de génération des tableaux (\texttt{04\_generated.py}) à partir de la table des termes de la base de travail, contrôlée par la recherche indépendante décrite au chapitre~\ref{ch:typologie}.

Les formules \emph{in pace} ne sont pas reprises ici : leurs listes complètes figurent au chapitre~\ref{ch:formulaire}. Les signes figurés sont décrits au chapitre~\ref{ch:supports}.

{\small
% généré par scripts/04_generated.py
\subsection*{Désignation de Dieu (19 occurrences)}\addcontentsline{toc}{subsection}{désignation de Dieu}
\begin{itemize}[leftmargin=1.2em,itemsep=0pt]
\item \nnum{214} (380) \lat{ancilla dei}, mot lu
\item \nnum{306} (389) \lat{ancilla dei}, mot lu
\item \nnum{458} (401) \lat{ancilla dei}, mot lu
\item \nnum{104} (366) \lat{dei}, partiellement restitué
\item \nnum{214} (380) \lat{dei}, mot lu
\item \nnum{306} (389) \lat{dei}, mot lu
\item \nnum{382} (395) \lat{dei}, entièrement restitué
\item \nnum{458} (401) \lat{dei}, mot lu
\item \nnum{474} (403) \lat{dei}, mot lu
\item \nnum{73} (363) \lat{deo}, mot lu
\item \nnum{73} (363) \lat{deum}, partiellement restitué — \textit{restitution de l’EDCS ; l’EDB imprime [quaeque de]metuisti}
\item \nnum{475} (403) \lat{deus}, mot lu
\item \nnum{475} (403) \lat{deus}, mot lu
\item \nnum{400} (397) \lat{domini}, mot lu
\item \nnum{85} (364) \lat{domino}, mot lu — \textit{lettres soulignées par l’EDB : lues par des éditeurs antérieurs, non visibles aujourd’hui}
\item \nnum{474} (403) \lat{homo dei}, mot lu
\item \nnum{232} (382) \lat{in domino}, développé d’une abréviation — \textit{développement assorti d’un point d’interrogation par l’éditeur}
\item \nnum{447} (401) \lat{θεὸν}, entièrement restitué
\item \nnum{447} (401) \lat{παντοκράτορα}, partiellement restitué
\end{itemize}
\subsection*{Désignation du Christ (9 occurrences)}\addcontentsline{toc}{subsection}{désignation du Christ}
\begin{itemize}[leftmargin=1.2em,itemsep=0pt]
\item \nnum{38} (352) \lat{Christi}, signe (S) — \textit{monogramme décrit par l’EDB et par l’estampage ILCV 3252A ; non compté comme mot}
\item \nnum{43} (355) \lat{Christi}, entièrement restitué
\item \nnum{140} (370) \lat{Christi}, développé d’une abréviation — \textit{chrisme décrit par l’EDB : monogramme probable}
\item \nnum{151} (371) \lat{Christi}, développé d’une abréviation — \textit{chrisme décrit par l’EDB : monogramme probable}
\item \nnum{306} (389) \lat{Christi}, partiellement restitué
\item \nnum{306} (389) \lat{Christi}, mot lu
\item \nnum{458} (401) \lat{Christi}, développé d’une abréviation
\item \nnum{311} (388) \lat{Christo}, développé d’une abréviation — \textit{chrisme décrit par l’EDB : monogramme probable}
\item \nnum{352} (393) \lat{Christum}, partiellement restitué
\end{itemize}
\subsection*{Esprit saint (2 occurrences)}\addcontentsline{toc}{subsection}{Esprit saint}
\begin{itemize}[leftmargin=1.2em,itemsep=0pt]
\item \nnum{249} (383) \lat{sanctum}, partiellement restitué — \textit{l’identification de l’Esprit saint dépend de cette restitution}
\item \nnum{249} (383) \lat{spiritum}, mot lu
\end{itemize}
\subsection*{Mention de la Trinité (1 occurrence)}\addcontentsline{toc}{subsection}{mention de la Trinité}
\begin{itemize}[leftmargin=1.2em,itemsep=0pt]
\item \nnum{474} (403) \lat{trinitatem}, mot lu
\end{itemize}
\subsection*{Initiation et foi (20 occurrences)}\addcontentsline{toc}{subsection}{initiation et foi}
\begin{itemize}[leftmargin=1.2em,itemsep=0pt]
\item \nnum{405} (397) \lat{catechumino}, mot lu
\item \nnum{70} (362) \lat{fide}, mot lu
\item \nnum{133} (369) \lat{fidei}, mot lu
\item \nnum{306} (389) \lat{fidei}, mot lu
\item \nnum{391} (396) \lat{fidei}, partiellement restitué
\item \nnum{70} (362) \lat{fidelis}, mot lu
\item \nnum{115} (367) \lat{fidelis}, mot lu
\item \nnum{231} (382) \lat{fidelis}, mot lu
\item \nnum{278} (383) \lat{fidelis}, mot lu
\item \nnum{391} (396) \lat{fidelis}, mot lu
\item \nnum{306} (389) \lat{fidem}, mot lu
\item \nnum{306} (389) \lat{fidem}, mot lu
\item \nnum{382} (395) \lat{fides}, mot lu
\item \nnum{172} (374) \lat{neofyta}, mot lu
\item \nnum{145} (371) \lat{neofyto}, mot lu
\item \nnum{210} (379) \lat{neofyto}, mot lu
\item \nnum{233} (382) \lat{neofyto}, mot lu
\item \nnum{276} (385) \lat{neofytus}, mot lu
\item \nnum{308} (389) \lat{neofytus}, mot lu
\item \nnum{391} (396) \lat{neofytus}, mot lu
\end{itemize}
\subsection*{Vocabulaire ecclésial (3 occurrences)}\addcontentsline{toc}{subsection}{vocabulaire ecclésial}
\begin{itemize}[leftmargin=1.2em,itemsep=0pt]
\item \nnum{18} (346) \lat{ecclesiae}, entièrement restitué
\item \nnum{67} (362) \lat{ecclesiae}, partiellement restitué
\item \nnum{428} (398) \lat{ecclesiae}, entièrement restitué
\end{itemize}
\subsection*{Église (épithète) (2 occurrences)}\addcontentsline{toc}{subsection}{Église (épithète)}
\begin{itemize}[leftmargin=1.2em,itemsep=0pt]
\item \nnum{67} (362) \lat{sanctae}, entièrement restitué
\item \nnum{428} (398) \lat{sanctus}, entièrement restitué
\end{itemize}
\subsection*{Saints (communion des saints) (4 occurrences)}\addcontentsline{toc}{subsection}{saints (communion des saints)}
\begin{itemize}[leftmargin=1.2em,itemsep=0pt]
\item \nnum{73} (363) \lat{sanctis}, mot lu — \textit{verbe restitué différemment : [vivit cum] (EDCS), [est sociatu]s (EDB)}
\item \nnum{330} (391) \lat{sanctis}, développé d’une abréviation — \textit{sigle p. t. c. s. : développement éditorial de quatre initiales}
\item \nnum{232} (382) \lat{sanctorum}, mot lu
\item \nnum{348} (392) \lat{ἁγίων}, mot lu
\end{itemize}
\subsection*{Intercession (2 occurrences)}\addcontentsline{toc}{subsection}{intercession}
\begin{itemize}[leftmargin=1.2em,itemsep=0pt]
\item \nnum{214} (380) \lat{ora}, mot lu
\item \nnum{457} (401) \lat{orationibus}, partiellement restitué
\end{itemize}
\subsection*{Résurrection (1 occurrence)}\addcontentsline{toc}{subsection}{résurrection}
\begin{itemize}[leftmargin=1.2em,itemsep=0pt]
\item \nnum{447} (401) \lat{ἀναστάσεως}, entièrement restitué — \textit{restitution contestée : l’EDB restitue κρίσεως}
\end{itemize}
\subsection*{Repos et sommeil de la mort (45 occurrences)}\addcontentsline{toc}{subsection}{repos et sommeil de la mort}
\begin{itemize}[leftmargin=1.2em,itemsep=0pt]
\item \nnum{179} (375) \lat{dormientes}, partiellement restitué
\item \nnum{73} (363) \lat{dormit}, mot lu
\item \nnum{400} (397) \lat{dormit}, mot lu
\item \nnum{464} (402) \lat{dormit}, partiellement restitué
\item \nnum{14} (345) \lat{Dormitio}, partiellement restitué
\item \nnum{134} (369) \lat{Dormitio}, mot lu
\item \nnum{134} (369) \lat{dormitio}, mot lu
\item \nnum{382} (395) \lat{pausat}, entièrement restitué
\item \nnum{458} (401) \lat{pausavit}, mot lu
\item \nnum{464} (402) \lat{pausavit}, partiellement restitué
\item \nnum{446} (400) \lat{quiescas}, développé d’une abréviation
\item \nnum{405} (397) \lat{quiescenti}, mot lu
\item \nnum{125} (368) \lat{Quiesces}, mot lu
\item \nnum{42} (355) \lat{quiescet}, mot lu
\item \nnum{59} (361) \lat{quiescet}, mot lu
\item \nnum{102} (366) \lat{quiescet}, mot lu
\item \nnum{73} (363) \lat{quiescis}, mot lu
\item \nnum{73} (363) \lat{quiescis}, mot lu
\item \nnum{231} (382) \lat{quiescis}, mot lu
\item \nnum{22} (348) \lat{quiescit}, mot lu
\item \nnum{70} (362) \lat{quiescit}, mot lu
\item \nnum{78} (363) \lat{quiescit}, partiellement restitué
\item \nnum{203} (377) \lat{quiescit}, mot lu
\item \nnum{204} (377) \lat{quiescit}, mot lu
\item \nnum{213} (380) \lat{quiescit}, développé d’une abréviation
\item \nnum{214} (380) \lat{quiescit}, mot lu
\item \nnum{370} (395) \lat{quiescit}, partiellement restitué
\item \nnum{371} (395) \lat{quiescit}, mot lu
\item \nnum{378} (395) \lat{quiescit}, mot lu
\item \nnum{391} (396) \lat{quiescit}, mot lu
\item \nnum{404} (397) \lat{quiescit}, mot lu
\item \nnum{421} (398) \lat{quiescit}, partiellement restitué
\item \nnum{428} (398) \lat{quiescit}, partiellement restitué
\item \nnum{450} (401) \lat{quiescit}, mot lu
\item \nnum{333} (391) \lat{requiescet}, mot lu
\item \nnum{474} (403) \lat{requiescet}, mot lu
\item \nnum{10} (344) \lat{requiescit}, partiellement restitué
\item \nnum{145} (371) \lat{requiescit}, mot lu
\item \nnum{374} (395) \lat{requiescit}, mot lu
\item \nnum{388} (396) \lat{requiescit}, mot lu
\item \nnum{429} (399) \lat{requiescit}, entièrement restitué
\item \nnum{460} (402) \lat{requiescit}, mot lu
\item \nnum{464} (402) \lat{requiescit}, partiellement restitué
\item \nnum{467} (402) \lat{requiescit}, mot lu
\item \nnum{258} (383) \lat{requiescunt}, mot lu
\end{itemize}
\subsection*{Allusion interprétative (niveau 5) (15 occurrences)}\addcontentsline{toc}{subsection}{allusion interprétative (niveau 5)}
\begin{itemize}[leftmargin=1.2em,itemsep=0pt]
\item \nnum{244} (383) \lat{aetheriam}, mot lu
\item \nnum{433} (399) \lat{astra}, entièrement restitué
\item \nnum{433} (399) \lat{astris}, entièrement restitué
\item \nnum{399} (397) \lat{Caelestinae}, mot lu
\item \nnum{382} (395) \lat{caelestis}, mot lu
\item \nnum{244} (383) \lat{caeli}, mot lu
\item \nnum{204} (377) \lat{Caelia}, mot lu
\item \nnum{73} (363) \lat{caelo}, entièrement restitué — \textit{restitution de l’EDCS ; l’EDB restitue [corpore namque tu]us}
\item \nnum{382} (395) \lat{caelo}, mot lu
\item \nnum{309} (389) \lat{luce}, mot lu
\item \nnum{400} (397) \lat{luce}, mot lu
\item \nnum{352} (393) \lat{lucis}, mot lu
\item \nnum{73} (363) \lat{lumine}, mot lu
\item \nnum{152} (371) \lat{redemptum}, partiellement restitué — \textit{lecture contestée par l’EDB}
\item \nnum{152} (371) \lat{redemptus}, mot lu — \textit{lecture contestée par l’EDB : [cor]p⟨o⟩re demptus}
\end{itemize}
\subsection*{Identité chrétienne (christianus) (1 occurrence)}\addcontentsline{toc}{subsection}{identité chrétienne (christianus)}
\begin{itemize}[leftmargin=1.2em,itemsep=0pt]
\item \nnum{201} (377) \lat{Christianae}, mot lu
\end{itemize}
\subsection*{Âme humaine (non doctrinal) (2 occurrences)}\addcontentsline{toc}{subsection}{âme humaine (non doctrinal)}
\begin{itemize}[leftmargin=1.2em,itemsep=0pt]
\item \nnum{73} (363) \lat{spiritus}, mot lu
\item \nnum{400} (397) \lat{spiritus}, mot lu
\end{itemize}
\subsection*{Épithète du défunt (non doctrinal) (5 occurrences)}\addcontentsline{toc}{subsection}{épithète du défunt (non doctrinal)}
\begin{itemize}[leftmargin=1.2em,itemsep=0pt]
\item \nnum{133} (369) \lat{sanctae}, mot lu
\item \nnum{399} (397) \lat{Sanctae}, mot lu
\item \nnum{310} (389) \lat{sanctissimae}, partiellement restitué
\item \nnum{26} (348) \lat{sanctitatis}, mot lu
\item \nnum{211} (379) \lat{sancto}, mot lu
\end{itemize}
\subsection*{Anthroponyme (niveau 6) (7 occurrences)}\addcontentsline{toc}{subsection}{anthroponyme (niveau 6)}
\begin{itemize}[leftmargin=1.2em,itemsep=0pt]
\item \nnum{297} (388) \lat{adeodatae}, anthroponyme
\item \nnum{112} (367) \lat{quodvuldeus}, anthroponyme
\item \nnum{388} (396) \lat{quodvultdeus}, anthroponyme
\item \nnum{80} (363) \lat{theodora}, anthroponyme
\item \nnum{213} (380) \lat{theodora}, anthroponyme
\item \nnum{5} (342) \lat{theodotus}, anthroponyme
\item \nnum{110} (367) \lat{theofila}, anthroponyme
\end{itemize}
\subsection*{Dominus (référent indéterminé) (1 occurrence)}\addcontentsline{toc}{subsection}{Dominus (référent indéterminé)}
\begin{itemize}[leftmargin=1.2em,itemsep=0pt]
\item \nnum{101} (366) \lat{domini}, partiellement restitué
\end{itemize}

}
